\skyrius{Problemos ir esamų sprendimų analizė}
\subsection{Poreikis ir informacijos šaltiniai}
\gaires{Analizuokite, ką patikimi šaltiniai padeda nuspręsti kuriant jūsų produktą.
Technologijų dokumentacija, moksliniai darbai ir naudotojų poreikių duomenys gali
atlikti skirtingą vaidmenį. Paaiškinkite, kodėl juos pasirinkote. Neužtenka
perpasakoti kelių tinklalapių.}
\pildyti{Šaltinių palyginimas, jų patikimumo ir tinkamumo įvertinimas bei
išvados produkto kūrimui. Jei rinkote poreikių duomenis, paaiškinkite iš ko,
kaip ir kokiomis sąlygomis juos gavote.}

\gaires{Citavimo pavyzdys su tikru šaltiniu: produkto kūrimo apraše galima pagrįsti
produkto poreikį situacijos analize \autocite{nsa-produkto-kurimas}.
Nuoroda į konkrečią vietą: \autocite[1~priedas]{nsa-instrukcija}.
Šie pavyzdžiai skirti parodyti LaTeX veikimą; pakeiskite juos savo temos šaltiniais.
Tiesioginę citatą rašykite kabutėse, o perfrazuotą mintį taip pat susiekite su šaltiniu.}

\subsection{Analogų palyginimas}
\gaires{Pasirinkite projektui reikšmingus palyginimo kriterijus. Įvardykite
sprendimų versijas ir šaltinius, o tada paaiškinkite, kuriuos jų sprendimus
perimsite, keisite ar atmesite. Lentelės eilučių skaičius nėra reikalavimas.}
\begin{table}[H]
\caption{Esamų sprendimų palyginimas}
\label{tab:analogai}
\begin{tabularx}{\textwidth}{@{}YYY@{}}
\toprule
Sprendimas ir šaltinis & Pranašumas pagal pasirinktą kriterijų & Trūkumas jūsų situacijoje \\
\midrule
\pildyti{Analogas A} & \pildyti{Faktas} & \pildyti{Faktas} \\
\pildyti{Analogas B} & \pildyti{Faktas} & \pildyti{Faktas} \\
\bottomrule
\end{tabularx}
\saltinis{\pildyti{autoriaus sudaryta pagal nurodytus šaltinius}}
\end{table}
\pildyti{Ką rodo \ref{tab:analogai} lentelė ir kokią konkrečią nišą užpildys
jūsų produktas? Venkite teiginio „nieko panašaus nėra“, jei jo negalite pagrįsti.}
